%% sections/method.tex -- Phase 7.1. Register: CLAUDE.md S9.2 (no em/en dashes as prose
%% punctuation). Encodes: PT-CS stratification = curation; the hybrid reporting rule (stated
%% ONCE, here); the S5.0 conventions; the MANDATORY QWK calibration sentence (first use of QWK
%% in the paper); SemEval fixed-0.5 + pooled accuracy/AUROC (per-split CUT 2026-06-15, CLAUDE.md S11);
%% backend reporting.
%% Every number here is sourced from the engine, matrix.yaml, runs.jsonl, or the generated tables.

\section{Study Design and Method}
\label{sec:method}

This section describes how the study was run. We set out the configuration design space
(Section~\ref{sec:design}), the four
datasets and the reliability of their reference grades (Section~\ref{sec:datasets}), the models and
serving backends (Section~\ref{sec:models}), the grading harness that drives every model through one
interface (Section~\ref{sec:harness}), and the metrics and statistical protocol applied throughout
(Section~\ref{sec:metrics}). Grading and scoring are identical across every model and configuration,
so that the comparisons the study draws are like-for-like.

\subsection{Design Space}
\label{sec:design}

The study crosses six configuration dimensions under one grading protocol
(Table~\ref{tab:design}). \emph{Reasoning} switches the model's thinking mode off or on.
\emph{Evaluation guidance} controls what grading material accompanies the question: the question stem
and the student answer are always sent, and the guidance level adds either nothing or the rubric (the
scoring criteria) or reference answer the dataset provides. \emph{Scope} contrasts grading a whole multi-question exam in one call
against one call per question. \emph{Decomposition} contrasts one holistic score against scoring each
rubric criterion separately and summing. \emph{Model family} covers three open-weight families and one
closed frontier model. \emph{Conversation state} contrasts a clean conversation per answer against a
single conversation that accumulates the grading history, and is run as a focused sub-study rather
than a factorial dimension. Answer order is controlled (natural versus inverse) because shared history
makes position a potential confound.

\begin{table}[!t]
\centering
\footnotesize
\caption{The configuration design space and the datasets each dimension applies to.}
\label{tab:design}
\begin{tabular}{lll}
\hline
Dimension & Levels & Applies to \\
\hline
Model family & 4 families (\S\ref{sec:models}) & all datasets \\
Reasoning & off, on & all datasets \\
Evaluation guidance & none, with guidance & all datasets \\
Scope & q-by-q, whole exam & PT-CS only \\
Decomposition & holistic, criterion & PT-CS only \\
Conversation state & clean, shared history & PT-CS sub-study \\
\hline
\end{tabular}
\end{table}

Two applicability rules keep the design honest. First, scope and decomposition exist only where the
data support them: PT-CS is the only dataset with multi-question submissions and per-criterion
rubrics, so those two dimensions run on PT-CS alone. Second, guidance means different material in the
two domains: a question-specific rubric for the code datasets (PT-CS, RIAYN) and a reference answer
for the short-answer datasets (Mohler, SemEval). The guidance effect is therefore two related
interventions, not one, and we never pool the two. PT-CS has no reference
solutions, and we did not synthesise any, to avoid introducing a confound.

\subsection{Datasets and Ground Truth}
\label{sec:datasets}

We develop every comparison on three public datasets and reserve PT-CS, our deployment dataset, for
transfer validation and for the experiments only it can support, which need its multi-question
submissions and rubric criteria.
\emph{Mohler}~\cite{mohler_text_2009,mohler_learning_2011} provides 2{,}273 CS short answers over 80
questions, each graded 0 to 5 by two teachers (we use their mean). The two annotators themselves agree
only moderately~\cite{mohler_learning_2011}, so agreement against their mean carries a human ceiling
that we keep in mind when reading absolute Mohler scores. \emph{SemEval-2013
Task~7}~\cite{dzikovska_semeval_2013} (SciEntsBank and Beetle) provides short answers with binary
correct/incorrect gold labels (the human reference a grader is scored against, which we call the gold).
SemEval is far larger than our other corpora, so it is the one we subsample: we draw 800 items by a
fixed-seed sample that preserves the dataset's question and domain proportions. The other public
datasets are graded in full. \emph{Rubric Is All You Need}
(RIAYN)~\cite{pathak_rubric_2025} provides 149 Java object-oriented-programming and data-structures
submissions with question-specific rubrics, graded by two graders in consensus; it is our public code
comparator. Published results situate our numbers but are not the success criterion
(Table~\ref{tab:published-baselines}): they come from different models, prompts, and protocols, so
all conclusions here rest on our own like-for-like contrasts.

\emph{PT-CS} is an export of a real Portuguese university grading deployment: 737 Java code answers
and 437 theory answers, in Portuguese, with question-specific rubrics (no reference solutions) and
heterogeneous point scales (1 to 12, normalised before any analysis). The export was pseudonymised at
ingest: all direct identifiers were dropped and every platform-generated comment was stripped, so no
model output from the platform reaches the graders we test. The platform seeds each grade with an artificial-intelligence (AI)
suggestion that teachers can review and correct, and the final grade is a two-teacher consensus with
no per-teacher records, so no inter-rater agreement is recoverable.

The reliability of the PT-CS gold is limited and uneven, and we treat that as a design input rather
than a footnote. Not all grades were teacher-reviewed; some may be unreviewed AI output. We therefore
stratify the 1{,}184 graded responses by an auditable criterion: a response is \emph{intervened} when
its final grade differs from the sum of its per-criterion scores, which is direct evidence that a human
moved the grade. Table~\ref{tab:strata} reports the strata: 382 responses (32.3\%) are intervened, 393
(33.2\%) match the sum exactly (suspected unreviewed), and 409 (34.5\%) have no per-criterion record at
all. We promote the intervened stratum to a curated transfer set in its own right,
\emph{PT-CS-verified} (295 code and 84 short-answer items in the analysed cells), and exclude the rest
from transfer claims by this declared criterion. The criterion detects intervention, not correctness:
an exact-sum grade may still be a teacher who reviewed and agreed, so the stratum is a conservative
proxy for review evidence, not a validation guarantee.

% auto-generated by make_phase5_tables.py -- do not hand-edit
\begin{table}[!t]\centering\small
\caption{PT-CS gold reliability strata. The intervened stratum, the only one with evidence of human review, is PT-CS-verified; full PT-CS is all three strata (1,184 responses).}\label{tab:strata}
\begin{tabular}{lll}
\hline
Stratum & n & \% \\
\hline
no-criteria (unassessable) & 409 & 34.5 \\
exact-sum (suspected unreviewed) & 393 & 33.2 \\
intervened (PT-CS-verified) & 382 & 32.3 \\
\hline
\end{tabular}
\end{table}


Because the PT-CS gold is of uneven reliability, every PT-CS agreement result we report uses
PT-CS-verified, the teacher-reviewed (intervened) subset defined above; the unreliable strata are
excluded by that declared criterion. The conversation sub-study (Section~\ref{sec:rq5}) needs no
stratum: it measures model behaviour and does not depend on the reference grades at all.

\subsection{Models and Backends}
\label{sec:models}

Four model families grade every dataset. Three are open-weight and expose a within-family reasoning
toggle, so the reasoning contrast holds the model fixed and flips only its thinking mode:
\emph{Qwen3.5} (served as \texttt{Qwen/Qwen3.5-35B-A3B}), \emph{DeepSeek-V4-Flash}
(\texttt{deepseek-ai/DeepSeek-V4-Flash}), and \emph{GLM-5.1} (\texttt{zai-org/GLM-5.1}), all run via
the DeepInfra inference application programming interface (API) with \texttt{enable\_thinking} off or
on. The fourth, \emph{GPT-5.1} (via
OpenAI, \texttt{reasoning\_effort} mapped to \texttt{none} or \texttt{high}), is a closed frontier
anchor: a reference point run on reduced item sets, used to situate the open models, never as an
inference target. Qwen3.5 with reasoning off, with guidance, question-by-question, and holistic
scoring is the internal baseline configuration against which the dimension contrasts are framed.

All calls use temperature 0 with a 4{,}096-token output ceiling for reasoning-off and 32{,}768 for
reasoning-on cells; every run row logs the provider, served model identifier, sampling parameters,
prompt hash, and timestamp (main matrix executed June 2026). We log these because absolute scores are conditional on the full serving setup, not just the prompt
and sampling we control: the served quantization, the inference software, and the hardware differ
across providers and cannot be matched exactly, so a different setup can yield different absolute
scores. The transferable claim is the configuration guidance (which configuration beats which), not
the absolute scores.

\subsection{Grading Harness}
\label{sec:harness}

A single grading harness, the software that runs every model through one common interface and records
its output, drives the whole study. Prompt templates are keyed by domain,
guidance level, scope, and decomposition; all templates request a machine-parseable JSON (JavaScript Object Notation) score, one
score per question for whole-exam calls and a per-criterion breakdown for criterion-level calls. A
deterministic parser extracts the grade and records whether extraction succeeded. We define the
\emph{extractable rate} $\pi$ as the fraction of runs whose output yields a parseable grade;
unparseable output is logged and excluded, never coerced to a zero score, because silent coercion
fabricates strictness. Each measurement is repeated $k=5$ times per item and configuration ($k=3$ for
the closed anchor), and results are cached by item, configuration hash, and repetition index, so reruns
never re-bill completed calls. Whole-exam calls grade many questions in one
API call, so token and cost are counted once per call, not once per question.

The sampled matrix balances power and cost. Reasoning-off cells run on the full public and PT-CS
datasets (2{,}273 Mohler, 149 RIAYN, 737 PT-CS code, 437 PT-CS theory items) and on the 800-item
SemEval sample, all at $k=5$. Reasoning-on is the most expensive axis, because reasoning traces run long, so those cells use a fixed
sample of 175 items per dataset and domain. The whole-exam arm grades 60 PT-CS submissions
(252 questions), and the anchor grades 60-item subsets at $k=3$. The conversation sub-study runs 20
submissions in two fixed answer orders (natural and inverse) on two models (Qwen3.5 and GLM-5.1) at
$k=3$, with everything else frozen at the baseline configuration. We did not run every combination.
Model family and reasoning are fully crossed: each of the four families grades with reasoning off and
on. The other dimension contrasts (evaluation guidance, scope, and decomposition) are secondary to
the reasoning question, so each was varied on its own from the baseline configuration, on the baseline
model (Qwen3.5), rather than crossed with every family, which would multiply the matrix. We therefore
report each of these effects on its own, not how two of them combine, and we make no claims about
combined effects.

\subsection{Metrics and Statistical Protocol}
\label{sec:metrics}

Two conventions are applied to every score before any metric. Predicted scores are clamped to
$[0, \text{scale max}]$, because rubrics with penalty criteria can push a criterion sum below zero
(88 such rows) and a grade cannot be negative. Scores are then normalised to $[0,1]$ by their scale
maximum, because the native scales differ (Mohler 5, SemEval 1, PT-CS 1 to 12) and raw scores must
never be pooled.

Agreement is reported primarily as Quadratic Weighted Kappa (QWK), the standard measure for ordinal
scoring: it penalises large grader-reference disagreements more than small ones. It reaches 1 at
perfect agreement and 0 at chance, and can in principle fall below 0 when a grader disagrees with the
reference more than chance would; all the values we report are positive. QWK needs discrete, ordered grades, but the model
produces a continuous score, so we first round each normalised score onto a common 0-to-5 ordinal
scale ($K{=}5$), the resolution of Mohler's native grades. Using the same scale for every dataset
makes the QWK values comparable. We fixed $K$ in advance, and recomputing with $K{=}4$ or $K{=}6$
changes no baseline QWK by more than 0.033, so the choice does not drive the results. QWK is then
computed on the mean of each item's $k$ repetitions.

Alongside QWK we report mean absolute error (MAE) on the normalised score, the average grading error
as a fraction of the scale, and Spearman rank correlation. Consistency is the mean within-item
standard deviation across the $k$ repetitions (SD$_k$); it needs no reference, and a lower value means
more reproducible grading. Cost is measured in output tokens and reported as the reasoning on/off
premium within each provider, since DeepInfra and OpenAI account for reasoning tokens differently. We also
track the signed model-minus-gold deviation, an item-level leniency measure used in
Section~\ref{sec:threats}.

SemEval needs its own mapping, because the model emits a continuous score while its gold is binary
(correct or incorrect). We turn the score into a correct/incorrect label at a fixed 0.5 cutoff, set in
advance rather than tuned to the data, because choosing the cutoff that maximised accuracy here would
inflate the result. We then report pooled accuracy together with AUROC, the area under the ROC curve,
a threshold-free measure equal to the chance a correct answer outscores an incorrect one (0.5 is
chance, 1 is perfect); it confirms the result does not hinge on the 0.5 cut.
We do not break SemEval into its seen and unseen subsets: those splits were defined for supervised
systems relative to a training set, and our prompted models have none, so the labels do not denote
novelty to our grader.

Every expensive contrast is computed on its paired subset, stated where the result is reported: the
reasoning contrast pairs off and on cells on the 175 sampled items (restricted to those whose
reasoning never overflowed the token ceiling, Section~\ref{sec:rq1}), the scope contrast on the 252
whole-exam questions, and anchor comparisons on the anchor's 60-item subsets. Uncertainty is given by
bootstrap 95\% confidence intervals (CIs) over items (800 resamples, fixed seed), and we call a
contrast significant when its CI excludes zero. We report effect sizes with CIs rather than p-values,
and apply no multiple-comparison correction; instead of family-wise claims we enumerate which cells
reach significance. The conversation sub-study uses paired Wilcoxon signed-rank tests, the standard
nonparametric test for paired differences.
